\section{Temporal structure of the SOS Newton system}\label{sec:structure}

\paragraph{Overview of the method.}
MIST factors $H_\delta$ in four steps, each justified by one result.
(i) The sparsity of $H_\delta$ is contained in a graph read directly off the
coefficient-matching equations (\cref{lem:envelope}).
(ii) Cutting the horizon between two stages separates this graph by a set of
equations whose size depends only on the variables shared across the cut
(\cref{thm:monomial}); with bounded temporal support this already gives an
ordering with linear work (\cref{thm:interval}).
(iii) MIST keeps a sequence of such interfaces, eliminates everything between
them independently, and obtains an exact block-tridiagonal system on the
interfaces plus a small border; a linear-time certificate guarantees this
structure for the interfaces actually chosen (\cref{thm:chain}).
(iv) The reduced system is solved by dense Cholesky or block cyclic reduction,
and the cost of the whole scheme is bounded by the local elimination fronts and
interface widths (\cref{thm:main}). This section covers steps (i) and (ii),
\cref{sec:hierarchy} step (iii), and \cref{sec:complexity} step (iv).

\subsection{The resource graph}

For each Gram block $c$ and free variable $a$, define the equations they touch:
\begin{equation}\label{eq:support}
\mathcal R_c=\{i:A_i^{(c)}\ne0\},\qquad \mathcal R_a=\{i:E_{ia}\ne0\}.
\end{equation}
We call a Gram block or a free variable a \emph{resource}.

\begin{definition}[resource graph]\label{def:resource}
The resource graph $\mathcal G=(V,\mathcal E)$ has one vertex per equation, and
each resource support $\mathcal R_c$ or $\mathcal R_a$ induces a clique.
\end{definition}

\begin{lemma}[structural envelope]\label{lem:envelope}
The off-diagonal sparsity pattern of $H_\delta$ is contained in $\mathcal G$. For
any elimination order, the fill of $H_\delta$ is contained in the symbolic fill of
$\mathcal G$ for the same order.
\end{lemma}

\begin{proof}
If equations $i$ and $j$ share no Gram block, every term in~\eqref{eq:schur}
vanishes. The term $\eta^2EE^\top$ in~\eqref{eq:operator} couples $i$ and $j$ only
if they share a free variable, and the shift adds nothing off the diagonal. A
pivot updates only pairs of its later neighbors, so induction over the elimination
order gives the fill statement.
\end{proof}

The rank-one term $uu^\top$ is deliberately excluded: it would make $\mathcal G$
complete. Sharing a resource allows a nonzero entry without forcing one, and two
equations using disjoint coordinates of the same Gram block are generally coupled
because $W_c$ is dense. The analysis therefore uses the conservative envelope and
never relies on cancellation.

\subsection{A monomial separator bound}

\begin{assumption}[local coefficient representation]\label{ass:local}
Each equation is identified by a unique monomial $\xi^\alpha$ with
$|\alpha|\le2r$. Each resource is assigned to one stage, and every monomial in its
support uses only variables of that resource. Equations outside this
representation are placed in a border set $B$.
\end{assumption}

At a temporal cut $t$, split the resources into those assigned to stages up to
$t$ and those after $t$. Let $V_L(t)$ and $V_R(t)$ be the variables used by the two
groups, and let
\begin{equation}\label{eq:shared}
S_t=V_L(t)\cap V_R(t),\qquad s_t=|S_t|.
\end{equation}
Let $\mathcal Q_t$ be the equations outside $B$ that touch resources on both sides.

\begin{theorem}[monomial interface bound]\label{thm:monomial}
Under \cref{ass:local}, removing $\mathcal Q_t$ disconnects the equations used only
on the left from those used only on the right in $\mathcal G[V\setminus B]$, and
\begin{equation}\label{eq:monomial}
|\mathcal Q_t|\le\binom{s_t+2r}{2r}.
\end{equation}
\end{theorem}

\begin{proof}
An equation in $\mathcal Q_t$ is a monomial occurring in a resource on each side,
so its variables lie in $V_L(t)\cap V_R(t)=S_t$. There are exactly
$\binom{s_t+2r}{2r}$ monomials of degree at most $2r$ in $s_t$ variables, and each
equation is a distinct monomial. Any remaining equation touches resources on one
side only. An edge between a left-only and a right-only equation would mean a
shared resource, which lies on one side and contradicts the other equation being
exclusive to the opposite side.
\end{proof}

The bound depends on the variables shared across the cut, not on the size of a
stage: a stage with $n_{\mathrm{var}}$ variables can carry up to
$\binom{n_{\mathrm{var}}+2r}{2r}$ equations, but the interface it creates is
controlled by $\binom{s_t+2r}{2r}$, and term sparsity typically keeps fewer
monomials still. In PushT (\cref{sec:experiments}), consecutive stages share only
the planar pose $(\cos\theta,\sin\theta,s_x,s_y)$, so $s_t=4$ and $r=2$ give at
most $\binom{8}{4}=70$ interface equations per cut, although each stage uses 22
variables. The interfaces MIST retains at $N=100$ after pruning and coarsening
have 37 equations on average.

\subsection{Bounded temporal support and an explicit linear-work ordering}

For an equation $i\notin B$, let $\mathcal T_i$ be the stages of the resources it
touches, and define its occurrence interval
\begin{equation}\label{eq:interval}
J_i=[a_i,b_i],\qquad a_i=\min\mathcal T_i,\quad b_i=\max\mathcal T_i.
\end{equation}

\begin{assumption}[bounded temporal support]\label{ass:support}
There are constants $\ell,q,g$, independent of $N$, such that $|B|=g$, every
equation $i\notin B$ satisfies $b_i-a_i+1\le\ell$, and at most $q$ equations first
appear at any one stage.
\end{assumption}

For~\eqref{eq:pop}, a variable used by stages $t$ and $t+1$ only produces
monomials of span at most two, and a stage with $n_{\mathrm{var}}$ variables gives
$q\le\binom{n_{\mathrm{var}}+2r}{2r}$. Monomials in global variables alone, such
as the constant monomial, go to $B$.

\begin{theorem}[interval path decomposition]\label{thm:interval}
Under \cref{ass:support}, $\mathcal G$ has a path decomposition of width at most
$\ell q+g-1$. Eliminating the equations outside $B$ in nondecreasing order of
$b_i$, followed by $B$, costs
\begin{equation}\label{eq:interval-cost}
W_{\mathrm{chol}}=O\bigl(Nq(\ell q+g)^2+g^3\bigr),\qquad
M_{\mathrm{chol}}=O\bigl(Nq(\ell q+g)+g^2\bigr).
\end{equation}
\end{theorem}

\begin{proof}
For each stage $t$, form the bag $X_t=B\cup\{i:t\in J_i\}$. Two non-border
equations sharing a resource lie together in the bag of that resource's stage, and
border equations lie in every bag, so all edges are covered. The bags containing a
non-border equation form its interval $J_i$, which gives the running-intersection
property. If $t\in J_i$, then $a_i$ is one of the $\ell$ stages ending at $t$, so
$|X_t|\le\ell q+g$.

For the ordering, consider the supergraph in which two non-border equations are
adjacent whenever their intervals intersect and border equations are adjacent to
everything. When the remaining equation with the smallest right endpoint $b_i$ is
eliminated, all of its remaining neighbors contain $b_i$ in their intervals, so
they already form a clique with $B$. The order creates no fill in the supergraph,
hence the fill of $\mathcal G$ stays inside it. Each of the at most $Nq$
non-border pivots has a front of size at most $\ell q+g$, costing
$O((\ell q+g)^2)$ work and $O(\ell q+g)$ storage; the border costs $O(g^3)$ and
$O(g^2)$.
\end{proof}

\Cref{thm:interval} is the multistage analogue of the classical link between
treewidth and factorization cost, with the decomposition read off the SOS model;
as noted by \citet{zhang2024chordal}, it is the incidence of constraints on
blocks, not the aggregate sparsity of the Gram matrices, that controls the
Newton-system cost. The theorem guarantees its own explicit order, not the result
of a generic AMD or nested-dissection ordering. It is also pessimistic: it charges
every pivot the largest bag, even though most equations of a stage belong to small
substructures. The hierarchy of the next section removes this pessimism.

Separators at different cuts need not be disjoint, since a monomial can persist
across several cuts, and separators chosen independently at each cut need not
produce a block-tridiagonal reduced system. The next section therefore makes
compatibility an explicit, checkable condition on the interfaces the solver uses.
